%% notes-tex/031-unified-representation/sections/1-question.tex

\section{The question and the decisions\secstatus{tent}}
\label{sec:question}

Five chemogenomic datasets are served and together hold 7.48 million records. Training one
model on them is attractive for an obvious reason and blocked by a less obvious one. The
attraction is scale and chemical breadth: no single dataset has both a genome-wide gene panel
and a large compound panel. The block is that a model takes one input format and one output
format, and the five datasets disagree about what a strain is, what a dose is, and what a
response is.

Three axes carry the input disagreement. Ploidy runs from haploid single deletions through
homozygous diploids to heterozygous diploids where one wild-type copy survives. Compound panels
run from 41 to 5{,}170. Dose runs from an unquantified target inhibition, through a molar
range spanning nine orders of magnitude, to a single fixed concentration applied to every
compound. The output disagreement is a fourth axis: the response is a log ratio in two
datasets, a z score in two, and a sensitivity score in one, read from a pooled barcode assay
in four and from liquid optical density in one.

The application target is Vanacloig-Pedros 2022 (\texttt{10.1128/spectrum.02017-22}), the
only served screen that doses the industrial hydrolysate panel at a matched target inhibition
and the only one in which isobutanol appears as a stress. The others enter as partners:
Hillenmeyer 2008 (\texttt{10.1126/science.1150021}) and Hoepfner 2014
(\texttt{10.1016/j.micres.2013.11.004}) for record count and gene coverage, Wildenhain 2015
(\texttt{10.1016/j.cels.2015.12.003}) for compound diversity.

\notesec{What this document is}
It is the planning document for the modeling phase, in one place: what the earlier data
exploration established about ceilings and transfer (Section~\ref{sec:established}), what a
shared input can absorb on each axis (Sections~\ref{sec:genotype} to \ref{sec:chemistry}),
which molecular encoders carry compounds into the model and how far they reach
(Sections~\ref{sec:encoders} and \ref{sec:metabolites}), how the channels enter the cell
graph transformer and what one decoder over four response scales looks like
(Section~\ref{sec:model}), and the representation itself with its implementation steps
(Section~\ref{sec:representation}).

\notesec{Two decisions taken}
The measurements below left two choices open, and both are now settled. The Hillenmeyer
homozygous arm is dropped: its 1.09 million records bring one non-exact chemical neighbor
above 0.5 to the application panel, and its heterozygous sibling keeps the gene coverage. Four
datasets remain, 6.39 million records. And the per-gene functional dose scales the existing
cross-attention perturbation operator rather than entering as a separate feature, because
the operator is already a soft mechanism and a scale factor that equals one for a deletion
leaves every current dataset's model unchanged.

\notesec{Scope}
This phase handles two perturbation kinds only: the full deletion and the one-of-two
copy-number variant, which are the only kinds any of the four datasets contains. Alleles off
the dosage axis, temperature-sensitive, decreased-abundance-by-mRNA-perturbation (DAmP) and
CRISPR interference, are not represented here and are not needed by these datasets; the
representation is written so they can be added without changing what is built now.
